% ========== Lecture 1 ========== %

%<*lec1-q0>

A fair coin is tossed twice. You have decided to play a game whose
outcome is

- If the first flip is H, you win \$1,000, and
- if the first flip is T, you lose \$1,000.

a. Writing a particular outcome as the concatenation of single outcomes
  such as H or T. What is the sample space $\Omega$?
b. Write as A the set corresponding to the event that the first flip is
H. What are elements in A? What is the probability of A?
c. Denote as B the same as above, of the event that the second toss is
T. What are the elements in B? What is the probability of B?
d. What is the probability of the first is head and second being tail?
e. Draw $A$, $B$, and $\Omega$ in a Venn diagram.
f. What is your expected return?

%</lec1-q0>

%<*lec1-q1>

If two balanced dice are rolled, what is the sample space? What is the set
  corresponding to the event that both rolls will be the same?


%</lec1-q1>

%<*lec1-q2>

From a survey of 60 students (undergraduate or graduate) attending a
  university, it was found that 9 were living off campus, 36 were
  undergraduates, and 3 were undergraduates living off campus. Find the number
  of these students who were

a. Undergraduates, were living off campus, or both.
b. undergraduates living on  campus.
c. graduate students living on campus.


%</lec1-q2>

%<*lec1-q3>

Draw Venn diagrams to verify:

a. $A = (A \cap B) \cup (A \cap B^C)$.

b. If $B \subset A$, then $A = B \cup (A \cap B^C)$.

%</lec1-q3>

%<*lec1-q4>

Show that, for any $A, B_1,\dots,B_n \in
  \Omega$ with $\cup_{i=1}^n B_i = \Omega$, we have $A = (A \cap B_1) \cup
  \cdots \cup (A \cap B_n)$.

%</lec1-q4>

%<*lec1-q5>

Consider the following experiments:

a. You have three keys and exactly one of them opens the door that you want
  to open; at each attempt to open the door, you pick one of the keys completely
  at random and try to open the door (at each attempt you pick one of the three
  keys at random, regardless of the fact that you may already have tried some in
  previous attempts). You count the number of attempts needed to open the
  door. What is the sample space for this experiment?
b. You have an urn with 1 red ball, 1 blue ball and 1 white ball. Two
  times in a row, you draw a ball from the urn, look at its color, take a note of the
  color on a piece of paper, and put the ball back in the urn. What is the
  sample space for this experiment? Write out the set corresponding to the event
  "you observe at least 1 white ball".


%</lec1-q5>

%<*lec1-q6>

If four fair coins are tossed, what is the probability
  that all four faces will be the same? Also, what is the probability that you
  will have the same number of both faces?

%</lec1-q6>


%<*lec1-q7>

Suppose a family contains two children of different ages, and we are interested in the gender
of these children. Let F denote that a child is female and M that the child is male and let a
pair such as F M denote that the older child is female and the younger is male. There are four
points in the set $\Omega$ of possible observations:
$$\Omega = \{FF, FM, MF, MM\}.$$
Let $A$ denote the subset of possibilities containing no males; $B$, the subset containing two
males; and $C$, the subset containing at least one male.
List the elements of $A$, $B$, $C$, $A \cap B$, $A \cup B$, $A \cap C$,
$A \cup C$, $B \cap C$, $B \cup C$, and $C \cap B^C$.

%</lec1-q7>


%<*lec1-q8>
%</lec1-q8>

% ========== Lecture 2 ========== %

%<*lec2-q1>

Use the axioms of probability to show that for any $A \subset \Omega$ that we
  have $P(A^C) = 1 - P(A)$.

%</lec2-q1>

%<*lec2-q2>

If two balanced dice are rolled, what is the probability
  that the sum of the two numbers that appear will be odd?


%</lec2-q2>

%<*lec2-q3>

If the probability that student 1 will fail a certain statistics
  examination is 0.5, the probability that student 2 will fail the examination
  is 0.2, and the probability that both student 1 and student 2 will fail the
  examination is 0.1, what is the probability that at least one of these two
  students will fail the examination?

%</lec2-q3>

%<*lec2-q4>

Let $B \subset A$. Show that $P(A) = P(B) + P(A\cap B^C)$.

%</lec2-q4>

%<*lec2-q5>

A fair die is tossed 6 times, and the number on the uppermost face is
  recorded each time. What is the probability that the numbers recorded are in
  the first four rolls are 1, 2, 3, 4 in any order? (Hint: It is useful to
  invoke the counting rules we learned in class.)

%</lec2-q5>

%<*lec2-q6>

Show that $P(\cup_{i=1}^n E_i) \le \sum_{i=1}^n P(E_i)$, for any sequence
  of events $\{E_i, i,\cdots, n\}$. (Hint: use induction.)

%</lec2-q6>

%<*lec2-q7>

Take an event $B \subset \Omega$. Rigorously prove the following three using
  probability axioms (or properties derived from them).

a. $P(B^C) = 1 - P(B)$.
b. If event $A$ *implies* (In other words, if event $A$ has
  happened, $B$ has also happened.) event $B$, $B$ is more probable than $A$,
  i.e., $P(A) \le P(B)$.
c. $0 \le P(A) \le 1$

%</lec2-q7>

%<*lec2-q8>
%</lec2-q8>

%<*lec2-q9>
%</lec2-q9>

%<*lec2-q10>
%</lec2-q10>

% ========== Lecture 3 ========== %

%<*lec3-q1>

If six dice are rolled, what is the probability that each of the six different numbers will
  appear exactly once?


  (Hint: what are the number of outcomes relevant for the denominator? For the
  numerator?)

%</lec3-q1>


%<*lec3-q2>

10 children are to be divided into an two teams (blue team and red team) of 5
  children each. Each team plays in a *separate* league.

a. How many different divisions are possible?
b. What if the two teams play in the same league?

%</lec3-q2>

%<*lec3-q3>

You have a string of 10 letters. How many substrings of length 7 can
  you form based on the original string?

%</lec3-q3>

%<*lec3-q4>

A company is looking to hire 3 new-grad analysts and asked UCSC for
  potential candidates. UCSC advised that they hire 3 students from the Fall
  STAT 131 class. Because the positions must be filled as quickly as possible,
  they decided to skip the interview process and hire 3 of you completely at
  random. What is the probability that *you* will become a software engineer
  at this company?

%</lec3-q4>


%<*lec3-q5>

Counting and arranging.

a. In how many different ways can the five letters $a$, $b$, $c$, $d$, and
  $e$ be arranged?
b. When rolling five dice, what's the probability that all five numbers
  shown will be different?
c. If six dice are rolled, what is the probability that each of the six
  different numbers will appear exactly once?

%</lec3-q5>


%<*lec3-q6>


  Calculate the probabilities of each of these (similar-sounding) events:
- At least one 6 appears when 6 fair dice are rolled.
- At least two 6's appear when 12 fair dice are rolled.

  Which of the two events has a higher probability?

%</lec3-q6>


%<*lec3-q7>

The STAT 131 final exam has 5 multiple choice questions. Questions 1 to
  3 and question 5 each have **3** possible choices. Questions 4 has **5** possible
  choices. Suppose that you haven't studied at all for the exam and you will
  pick your final exam answers completely at random.

a. What is the probability that you will get a full score at the exam?
b. What is the probability that you will get a score higher than 70\%?

%</lec3-q7>


%<*lec3-q8>

How many ways there are to put $n$ indistinguishable balls into $k$
  distinguishable bins? (Hint: you can't simply use a formula from lecture.)


%</lec3-q8>


%<*lec3-q9>

Say we are to randomly distribute distribute $n$ indistinguishable cookies
  to $k$ distinct kids. Some kids may get zero cookies. How many ways to do
  this?


%</lec3-q9>


%<*lec3-q10>

Cards are dealt at random and without replacement from a standard 52 card
  deck. What is the probability that the second king is dealt on the fifth card?


%</lec3-q10>


%<*lec3-q11>

**Practice** In the scenario from the previous problem, what is the probability of
  these three events?

a. Each kid receives at least one cookie.
b. Each kid receives at least two cookies. (assume $n\ge 2k$)

%</lec3-q11>


% ========== Lecture 4 ========== %
%<*lec4-q1>

Every time a customer buys a product, they select either brand $A$ or
  brand $B$. For each purchase, there's a 1/4 chance they'll buy the same brand
  as they did last time, and a 3/4 chance they'll switch brands.  If on the
  initial buy there's an equal chance of picking either brand $A$ or $B$, what's
  the probability that the customer will buy brand $B$ for their first two
  purchases and then choose brand $A$ for their third and fourth purchases?

%</lec4-q1>


%<*lec4-q2>

Solve these, and justify your answers.

- If $A \subset B$ with $\mathbb{P}(B) >0$, $\mathbb{P}(A|B) = ________.$
- If $A$ and $B$ are disjoint events with $\mathbb{P}(B) >0$,  $\mathbb{P}(A|B) = ________.$
- If $\mathbb{P}(A)>0$, $\mathbb{P}(A|A) = ________$

%</lec4-q2>


%<*lec4-q3>

The probability that any child in a certain family will have curly hair is
  1/4, and this feature is inherited independently by different children in the
  family. If there are five children in the family and it is known that at least
  one of these children has curly hair..

a. What is the probability that at least three of the children have
  curly hair?
b. If it is known that the youngest child in the family has curly hair,
  what is the probability that at least three of the children have curly hair?
c. Explain why the answer in part (a) is different from the answer in
  (b).

%</lec4-q3>


%<*lec4-q4>

Let $A,B \in \Omega$ $P(B)>0$.

a. What is $P(A|\Omega)$?
b. What is $P(A|A)$?
c. Let $ B \in A$. What is $P(A|B)$?
d. Let $A \cap B = \emptyset$. What is $P(A|B)$?

%</lec4-q4>


%<*lec4-q5>

Consider the following events and the corresponding table of probabilities:


|      | $B$ | $B^C$ |
|------|-----|-------|
| $A$   | 1/4 | 1/12 |
| $A^C$ | 1/2 | 1/6  |


Are the events A and B independent? Justify your answer.

%</lec4-q5>


%<*lec4-q6>

Assume there is a box containing five coins, each with a distinct
  probability of landing heads when tossed. Probabilities are denoted as $p_i$
  for the $i$th coin (where $i$ ranges from 1 to 5). Given the specific
  probabilities, $p_1 = 0, p_2 = 1/4, p_3 = 1/2, p_4 = 3/4$, and $p_5 = 1$.


a. Say you toss a randomly selected coin. If a head is obtained, what
  is the posterior probability that the $i$th coin was selected ($i=1,...,5$)?
b. What would be the probability of getting another head if the same
  coin were tossed once more?
c. If the first toss of the chosen coin resulted in a tail, and the
  same coin were tossed again, what would be the probability of getting a head
  on the second toss?

%</lec4-q6>


%<*lec4-q7>

Two persons are taking turns throwing a ball at a target. Person 1 has a
  $1/3$ probability of hitting the target with each throw, while Person 2 has a
  $1/4$ probability. Assuming person 1 starts the game, determine the
  probability that the target will be hit for the first time when person 1 makes
  their 4th throw.


%</lec4-q7>


%<*lec4-q8>

Let us say that $A$ and $B$ are conditionally independent on
  $C$. Prove that $P(A | B \cap C) = P(A | C)$.


%</lec4-q8>


%<*lec4-q9>

If you sample $n=4$ cards with replacement from $r=3$ distinct cards, what
  is the probability of obtaining a complete set?


%</lec4-q9>


%<*lec4-q10>

[**Practice**] Guessing gender. (Fun fact! This is a popular mathematical paradox by
  Martin Gardner in *Scientific American* in 1959.)

a. Mr. Jones has two children. The older child is a girl. What is the
  probability that both children are girls?
b. Mr. Smith has two children. At least one of them is a girl. What is the
  probability that both children are girls?

%</lec4-q10>


% ========== Lecture 5 ========== %

%<*lec5-q1>

Suppose that a box contains seven red balls and three blue balls. If five
  balls are selected at random without replacement,

a. What is the expected number of red balls that will be obtained?
b. Determine the pmf of the number of red balls that will be obtained.

%</lec5-q1>

%<*lec5-q2>

Suppose that two balanced four-sided dice are rolled, and let $X$ denote
  the sum of the two numbers that appear, and let $Y$ denote the absolute value
  of the difference between the two numbers that appear.

a. Derive the probability mass function (pmf) of the random variable $X$.
b. Also derive the pmf of the random variable $Y$. What is the most
  probable value of $Y$?
c. Sketch the probability mass function for the above two random
  variables $X$ and $Y$.

%</lec5-q2>


%<*lec5-q3>

An airline knows that $5$ percent of the people making reservations on a
  certain flight will not show up. Consequently, their policy is to sell $52$
  tickets for a flight that can hold only $50$ passengers. What is the
  probability that there will be a seat available for every passenger who shows
  up? (Calculate a final numerical answer for full credit.)

%</lec5-q3>


%<*lec5-q4>

Fun with probability mass functions.

a. Suppose that a random variable X has the Bernoulli distribution with
  parameter $p= 0.7$. Derive and sketch the probability mass function. Mark
  the expectation as a vertical line.
b. Suppose that a random variable X has the Binomial distribution with
  parameters $n=2$ and $p= 0.7$. Derive and sketch the probability mass
  function. Mark the expectation as a vertical line.
c. Calculate the probability $\mathbb{P}(X \in (1.9,3])$ and
  $\mathbb{P}(X \in (2,3])$ based on (b).

%</lec5-q4>

%<*lec5-q5>

Suppose $X$ has a binomial distribution with parameter $n=6$ and some $p$.

a. Show that if $p=1/2$, $X=3$ is the most likely outcome.
b. Derive the value of $p$ (smaller than $1/2$) at which the most
  likely outcome *ceases* to be $X=3$.


  than $3$.

%</lec5-q5>

%<*lec5-q6>

If 10 percent of the balls in a certain box are red, and if 20 balls are
  selected from the box at random with replacement, what is the probability that
  more than three red balls will be obtained?


%</lec5-q6>


%<*lec5-q7>

Suppose that an airplane engine will fail, when in flight, with
  probability $1-p$ independently from engine to engine. Suppose that the
  airplane will make a successful flight if at least $50\%$ of its engines
  remain operative.

a. What is the probability mass function of the (random) number of
  operative engines of a given plane with $n=5$ engines? Derive and sketch it
  for $p=0.1$.
b. For what values of $p$ is a four-engine plane preferable to a two-engine
  plane?
c. If $p$ is revealed to be $0.1$, is it more probable for a two-engine or
  a three-engine plane to have at least one engine fail?

%</lec5-q7>

%<*lec5-q8>

Suppose $X$ is the random variable denoting the outcome of a fair dice
  roll. Also suppose a barista at a coffee shop is serving six customers
  numbered $1,\cdots, 6$ and selects one at random (to gift a free muffin),
  calling that customer number $Y$.

a. Derive the probability mass functions of the two random variables
  $X$ and $Y$. (Hint: it is possible for two random variables to have the same
  *distribution*.)
b. Prove that these two random variables are *independent*, by
  proving that $\mathbb{P}(X=x, Y=y) = \mathbb{P}(X=x) \cdot \mathbb{P}(Y=y)$ for any $x,y$ values
  between $1$ and $6$.
b. What is the probability that $X=Y$?


%</lec5-q8>

%<*lec5-q9>

[**Practice**] A fair coin is independently flipped $10$ times, $3$
  times by person A, and $7$ times by person B.  Show that the probability that
  person A and person B flip the same number of heads is equal to the
  probability that there are a total of $3$ heads.


  (Hint 1: the event that person
  A gets $j$ heads is independent of the event that person B gets $j$
  heads. Then, what is the probability that person both A and B get $j$ heads?
  Hint2: the Vandermonde's identity could be useful here.)


  (More generally, prove that this holds for any fair coin flipped $n$ times,
  $k$ times by person A, and $n-k$ times by person B, where the above problem is
  the special case of $k=3$.)

%</lec5-q9>

%<*lec5-q10>

If a fair coin is successively flipped, find the probability that a head
   first appears on the fifth trial.

%</lec5-q10>

% ========== Lecture 6 ========== %
%<*lec6-q1>

Say you are a farmer waiting for it to rain two days in a row -- so you
  can sow your seeds -- where the daily probability of raining is $p$. What is
  the probability that you will have to wait four or more days in total?


%</lec6-q1>

%<*lec6-q2>

Suppose that a random variable $X$ can take only the two values $a$ and
  $b$ with the following probabilities: $\mathbb{P}(X= a)= p$ and $\mathbb{P}(X= b)= 1-p$.
  Express the pmf of $X$ in a form similar to that of the Bernoulli pmf in the
  lecture notes.

%</lec6-q2>


%<*lec6-q3>

A parking lot has two entrances. Cars arrive at entrance I according to a
  Poisson distribution at an average of three per hour and at entrance II
  according to a Poisson distribution at an average of four per hour.  What is
  the probability that a total of three cars will arrive at the parking lot in a
  given hour? Assume that the numbers of cars arriving at the two entrances are
  independent.

%</lec6-q3>


%<*lec6-q4>

Say you are waiting for the first day it will rain, where the daily
  probability of raining is $p$. Then the total number of days it will take to
  rain, $X$, follows a geometric distribution with parameter $p$.  Prove that
  this is the pmf of $X$:
  $$f(x \mid p)=
  \begin{cases}p(1-p)^x & \text { for } x=0,1,2, \ldots \\ 0 & \text { otherwise
    }\end{cases}$$

%</lec6-q4>


%<*lec6-q5>


  Suppose that a box contains $A$ red balls and $B$ blue balls. Suppose also
  that $n \ge 0$ balls are selected at random from the box without replacement,
  and let $X$ denote the number of red balls that are obtained. (Clearly, we
  must have $n \le A + B$ or we would run out of balls. Also, if $n = 0$, then
  $X= 0$ because there are no balls, red or blue, drawn.) If $n \ge 1$, we let
  the random variable $X_i$ be defined so that $X_i = 1$ if the $i$-th ball
  drawn is red and $X_i = 0$ if not.

a. What are $\mathbb{P}(X_2 = 1|X_1 = 1)$ and $\mathbb{P}(X_2 = 1|X_1 = 0)$?
b. True or False? $X_1, \cdots, X_n$ are mutually independent. Justify your
  answer.
c. Use combinatorial arguments to derive the pmf of $X$.


  (Make sure to precisely define the support of $X$.)

%</lec6-q5>


%<*lec6-q6>

At the US Postal office, the expected number of customers between 1pm and
  2pm is usually $3$. The number behaves like a Poisson random variable. What is
  the probability that 2 or more people will visit the postal office during that
  time period?


%</lec6-q6>


%<*lec6-q7>

Denote as $f_1(x \mid n,p)$ the pmf of a Binomial distribution with
  parameters $n$ and $p$. Also denote as $f_2(x \mid \lambda)$ the pmf of a
  Poisson distribution with parameter $\lambda$.
  Let $\{p_n\}_{n=1, \cdots, \infty}$ be a sequence of numbers between 0 and 1
  such that $\lim_{n \to \infty}np_n = \lambda$. Prove that:
  $$ \lim_{n\to\infty} f_1(x \mid n, p_n) = f_2(x \mid \lambda).$$
  Note, this *formally* justifies why we can use a Poisson distribution with
  mean $np$ to approximate a Binomial distribution with parameters $n$ and $p$
  where $n$ is large and $p$ is close to 0.

%</lec6-q7>


%<*lec6-q8>

Suppose that in a large population the proportion of people who have a
  certain rare disease is $0.01$. We shall determine the probability that in a
  random group of $200$ people at least four people will have the disease, in
  two ways:

a. Calculate the exact probability.
b. Calculate the approximate probability, using an appropriate Poisson
  distribution.

%</lec6-q8>


%<*lec6-q9>

Consider three urns. Urn A contains 2 white and 3 red balls, urn B
  contains 8 white and 4 red balls and urn C contains 1 white and and 3 red
  balls. If one ball is selected from each urn, what is the probability that the
  ball chosen from urn A was white, given that exactly 2 white balls were
  selected? (Hint: consider defining random variables $X_i$ (for $i=1,2,3$) that
  is valued 1 if a white ball is drawn from the $i$'th urn, zero otherwise.)

%</lec6-q9>

%<*lec6-q10>

If a random variable $X$ has a discrete distribution for which the pmf is
   $f(x)$, then the value of $x$ for which $f(x)$ is maximum is called the *mode*
   of the distribution. If this same maximum $f(x)$ is attained at more than one
   value of $x$, then all such values of $x$ are called modes of the
   distribution. Find the mode or modes of the binomial distribution with
   parameters $n$ and $p$. (Hint: Study the ratio
   $f(x+1 \mid n, p) / f(x \mid n, p)$.)

%</lec6-q10>


% ========== Lecture 7 ========== %
%<*lec7-q1>

Suppose that the pdf of a random variable $X$ is:
  $$f(x)=
  \begin{cases}
    \frac{1}{8} x &\text{ for }0 \le x \le 4 \\
    0 &\text{ otherwise }
  \end{cases}$$

a. Sketch the pdf.
b. Find the value of $t$ such that $\mathbb{P}(X \ge t)= 1/2$.
c. What is more probable -- a smaller (near zero) or larger value (near
  4) of $X$?

%</lec7-q1>


%<*lec7-q2>

Suppose that a random variable $X$ has a discrete distribution
   with the following probability mass function:

$$
f(x)=
     \begin{cases}
       cx & \text{for } $x = 0, 1, 2, 3, 4, 5.$\\
       0 & \text{otherwise}.
     \end{cases}
$$

   Determine the value of the constant $c$.

%</lec7-q2>


%<*lec7-q3>

Take random variable $X$ whose pdf is:
  f(x)= \begin{cases}c\left(9-x^2\right) & \text { for }-3 \leq x \leq 3, \\ 0 & \text { otherwise. }\end{cases}

a. Find the constant $c$ and sketch the pdf.
b. Calculate $\mathbb{P}(X<0)$, $\mathbb{P}(-1 \leq X \leq 1)$, and $\mathbb{P}(X>2)$.

%</lec7-q3>


%<*lec7-q4>

Take an exponential random variable $X$ with parameter $\beta$. What is the
  *cumulative* probability up till a point $x$, i.e., $\mathbb{P}(X \le x)$? Sketch
  and derive this.

%</lec7-q4>


%<*lec7-q5>

Suppose that a random variable $X$ has the uniform distribution on the
  interval $[-2, 8]$. Sketch the pdf of $X$ and the value of $\mathbb{P}(0 < X < 7)$.

%</lec7-q5>


%<*lec7-q6>

If a random variable $X$ has a distribution for which the pdf is $f(x)$,
  then the value of $x$ for which $f(x)$ is maximum is called the *mode* of
  the distribution. (If this same maximum $f(x)$ is attained at more than one
  value of $x$, then all such values of $x$ are called modes of the
  distribution.)

  Find the mode or modes of the gamma distribution with parameters $\alpha$ and
  $\beta$.

%</lec7-q6>


%<*lec7-q7>

Suppose that the pdf of a random variable $X$ is as follows:

$$
f(x) =
    \begin{cases}
      \frac{c}{(1-\frac{x}{2})^{\frac{1}{2}}} & \text{for }\ 0<x<1,\\
      0 & \text{otherwise}
    \end{cases}
$$


a. Find the value of the constant $c$ and sketch the pdf.
b. Find the value of $P(X\leq \frac{1}{2})$

%</lec7-q7>

%<*lec7-q8>

Evaluate the integral $\int e^{-3x^2} \;dx$. (Hint: see the pdf of a
  normal distribution).

%</lec7-q8>

%<*lec7-q9>

Sketch the beta pdf with the following parameters (hint: try calculating
  and plotting some $(x, f(x))$ points).

a. $\alpha=\beta=1$.
b. $\alpha=\beta=3$.
c. $\alpha=1$, $\beta=2$.

%</lec7-q9>


%<*lec7-q10>

The *expected value* of a random variable $X$ with pdf $f$ is defined as
  $\int_{\mathbb{R}} x \cdot f_X(x)\; dx$.  What is the expected value of
  a beta random variable with parameters $\alpha$, $\beta$?

%</lec7-q10>

%<*lec7-q11>

Let $X$ be a random variable with pdf $f$. The $p$-quantile function
  $g(p)$ of $X$ is defined as the smallest value $x$ of $X$ for which
  $\mathbb{P}(X \le x) \ge p$.  (For example, the $0.1$-th quantile is the smallest
  value $x$ of $X$ for which the probability to the left-hand-side of $x$,
  $P(X < x)$, is at least $0.1$.)  Compute the quantile function $g(p)$ of the
  exponential distribution with parameter $\beta$.

%</lec7-q11>

%<*lec7-q2>

Suppose that the proportion $X$ of defective items in a large lot is
   unknown and that $X$ has the beta distribution with parameters $\alpha$ and
   $\beta$. If one item is selected at random from the lot, what is the
   probability that it will be defective?

%</lec7-q12>

%<*lec7-q13>

Suppose that a certain test is to be taken by three students independently
  of one another, and the number of minutes required by any particular student
  to complete the test has the exponential distribution with parameter
  $1/40$. (This means that the *expected* value of the wait time is $40$
  minutes.)

a. Suppose that the test begins at 9:00 a.m. Determine the probability
  that at least one of the students will complete the test before 9:30 a.m.
b. Suppose again that the test is taken by three students, and the
  first student to complete the test finishes at 9:20 a.m. Determine
  the probability that at least one other student will complete the
  test before 10:00 a.m.
c. Suppose again that the test is taken by five students. Determine
  the probability that no two students will complete the test within 10
  minutes of each other.

%</lec7-q13>

%<*lec7-q14>

A manufacturer believes that an unknown proportion
 P of parts produced will be defective. She models P as
 having a beta distribution. The manufacturer thinks that P
 should be around 0.05, but if the first 10 observed products
 were all defective, the mean of P would rise from 0.05 to
 0.9. Find the beta distribution that has these properties.

%</lec7-q14>

%<*lec7-q15>

Suppose that X has the beta distribution with parameters α and β. Show
   that 1− X has the beta distribution with parameters β and α.

%</lec7-q15>


% ========== Lecture 8 ========== %
%<*lec8-q1>

Suppose that a fair coin is tossed repeatedly until a head is obtained for
  the first time, and let $X$ denote the number of tosses that are
  required.


a. What is probability $P(X=x)$ for $x=0,1,2,\cdots,$?
b. Sketch the cdf of $X$. (Hint: start by sketching the pmf).

%</lec8-q1>


%<*lec8-q2>

Suppose that the joint p.d.f. of X and Y has the follow-
 f (x, y)=
 kx2y2 for x2 + y2 ≤ 1,
 0 otherwise.
 Show (formally or informally) that X and Y are not independent.

%</lec8-q2>


%<*lec8-q3>

Take $X$, a random variable distributed as $\text{Unif}(0, 5)$. Also
   take a random variable $Y$ defined as follows: $Y=0$ if $X \le 1$, $Y= 5$ if
   $X\ge 3$, and $Y= X$ otherwise.  Sketch the cumulative distribution function
   (cdf) of $Y$.

%</lec8-q3>


%<*lec8-q4>

Draw both a pmf and cdf of the random variable $X\sim \text{Bern}(0.6)$.

%</lec8-q4>


%<*lec8-q5>

What is the median of the following two random variables?

a. $X\sim \text{Bern}(0.6)$.
b. $X\sim \text{Bern}(0.3)$.

%</lec8-q5>


%<*lec8-q6>

Suppose that the cdf of a random variable $X$ is as follows:


$$
F(x)= \begin{cases}e^{x-3} & \text { for } x \leq 3, \\ 1 & \text { for } x>3 .\end{cases}
$$


a. Find and sketch the pdf of $X$.
b. Find and sketch the cdf of $X$.
c. Find the first quartile.
d. Find the median.

%</lec8-q6>


%<*lec8-q7>

Suppose that $X$ has the uniform distribution on the interval $[0,5]$ and
  that the *mixed* random variable $Y$ is defined by $Y=0$ if
  $X \leq 1, Y=5$ if $X \geq 3$, and $Y=X$ otherwise. Sketch the cdf of $Y$.

  (Hint: what are probabilities $\mathbb{P}(Y=0)$, $\mathbb{P}(Y=5)$? What about $\mathbb{P}(Y=y)$
  for any other value $y \not\in \{0,5\}$? Use this to complete the cdf.)

%</lec8-q7>


%<*lec8-q8>

Suppose that a point in the $x y$-plane is chosen at random from the
  interior of a circle for which the equation is $x^2+y^2=1$; and suppose that
  the probability that the point will belong to each region inside the circle is
  proportional to the area of that region. Let $Z$ denote a random variable
  representing the distance from the center of the circle to the point. Find and
  sketch the cdf of $Z$.

%</lec8-q8>


%<*lec8-q9>

Suppose that the c.d.f. of a random variable $X$ is as follows:

$$
F(x)= \begin{cases}0 & \text { for } x \leq 0 \\ \frac{1}{9} x^2 & \text { for } 0<x \leq 3 \\ 1 & \text { for } x>3\end{cases}
$$


Find and sketch the p.d.f. of $X$.

%</lec8-q9>


%<*lec8-q10>

Prove that the quantile function $F^{-1}(\cdot)$ of a general random
  variable $X$ has the following property analogous to that of the cdf:
  $F^{-1}(p)$ is a non-decreasing function of $p$ for $0<p <1$.

%</lec8-q10>


%<*lec8-q11>

Take a random variable $X$ with an exponential distribution whose
  parameter is $\lambda$. What is the median?

%</lec8-q11>


%<*lec8-q12>

Suppose that the proportion $X$ of defective items in a large lot is
   unknown and that $X$ has the beta distribution with parameters $\alpha$ and
   $\beta$. If one item is selected at random from the lot, what is the
   probability that it will be defective?

%</lec8-q12>


% ========== Lecture 9 ========== %
%<*lec9-q1>

In a certain suburban area, each household reported the number of cars and
  the number of television sets that they owned. Let $X$ stand for the number of
  cars owned by a randomly selected household in this area. Let $Y$ stand for the
  number of television sets owned by that same randomly selected household. This
  is the joint pmf (each entry is $f(x,y)$):


| $X ackslash Y$ | 1 | 2 | 3 | 4 |
|------------------|---|---|---|---|
| 1 | 0.1 | 0 | 0.1 | 0 |
| 2 | 0.3 | 0 | 0.1 | 0.2 |
| 3 | 0 | 0.2 | 0 | 0 |


 What is the probability that a randomly selected household will

a. have more than two cars *and* more than two TVs?
b. have the same number of cars and TVs?
c. have exactly one car?

%</lec9-q1>


%<*lec9-q2>

Suppose that the joint p.d.f. of X and Y has the follow-
 f (x, y)=
 kx2y2 for x2 + y2 ≤ 1,
 0 otherwise.
 Show (formally or informally) that X and Y are not independent.

%</lec9-q2>


%<*lec9-q3>

Take $X$, a random variable distributed as $\text{Unif}(0, 5)$. Also
   take a random variable $Y$ defined as follows: $Y=0$ if $X \le 1$, $Y= 5$ if
   $X\ge 3$, and $Y= X$ otherwise.  Sketch the cumulative distribution function
   (cdf) of $Y$.

%</lec9-q3>


%<*lec9-q4>

The random variables X and Y have joint pdf $f_{X,Y}(x,y)= 0.5$ in the
  triangle of the x-y plane determined by the points $(0,0)$, $(1,0)$, and
  $(0,1)$. Call this triangle $T$.


a. Find the joint cdf of $(X,Y)$.


  (Hint 1 : if the pdf is constant, how do you calculate the probability
  $\mathbb{P}(X \le x, Y \le y)$? Hint 2: try sketching the integration region for
  various values of $(x,y)$.)


b. Verify that your joint cdf differentiates back to the pdf
  $f_{X,Y}(x,y)= 0.5 \cdot \mathbf{1}_{(x,y) \in T}$

%</lec9-q4>


%<*lec9-q5>

Suppose that the joint cdf of $X$ and $Y$ has the form
  $F(x,y) =\frac{1}{16}xy(x+y)$ when for values between $0$ and $2$.

a. Write out the *full* joint cdf -- for *all* values of $x \in \mathbb{R}$ and $y \in \mathbb{R}$?
b. What is the cdf of $X$?
c. What is the joint pdf of $X$ and $Y$?

%</lec9-q5>


%<*lec9-q6>

Suppose that $X$ and $Y$ are random variables such that ( $X, Y$ ) must
  belong to the rectangle in the $x y$-plane containing all points $(x, y)$ for
  which $0 \leq x \leq 3$ and $0 \leq y \leq 4$. Suppose also that the joint
  cdf of $X$ and $Y$ at every point $(x, y)$ in this rectangle is specified
  as follows:

$$
F(x, y)=\frac{1}{156} x y\left(x^2+y\right) .
$$


Determine


a. $\mathbb{P}(1 \leq X \leq 2$ and $1 \leq Y \leq 2)$.
b. $\mathbb{P}(2 \leq X \leq 4$ and $2 \leq Y \leq 4)$.
c. the cdf of $Y$.
d. the joint pdf of $X$ and $Y$.
e. $\mathbb{P}(Y \leq X)$.

%</lec9-q6>


%<*lec9-q7>

Consider an illuminated signboard where the top row has three bulbs and
  the bottom row contains four. Use \(X\) to denote the count of non-working
  bulbs in the top row at a specific time \(t\), and \(Y\) for the bottom row at
  that same time. The joint pmf of \(X\) and \(Y\) is shown in this table:


| $X ackslash Y$ | 0 | 1 | 2 | 3 | 4 |
|------------------|---|---|---|---|---|
| 0 | 0.08 | 0.07 | 0.05 | 0.02 | 0.01 |
| 1 | 0.06 | 0.10 | 0.12 | 0.05 | 0.02 |
| 2 | 0.05 | 0.06 | 0.11 | 0.02 | 0.03 |
| 3 | 0.04 | 0.01 | 0.03 | 0.03 | 0.04 |


	a. Calculate $P(X=1)$.
	b. Calculate $P(Y \geq 1)$.
	c. Calculate $P(X\leq 1 \text{ and } Y\leq 1)$.
	d. Calculate $P(X = Y)$.
	e. Calculate $P(X > Y)$.

%</lec9-q7>


%<*lec9-q8>

Let $Y$ be the rate (calls per hour) at which calls arrive at a
  switchboard. Let $X$ be the number of calls during a two-hour period. A
  popular choice of joint pmf/pdf for ($X, Y$) in this example would be
  one like
$$
f(x, y)= \begin{cases}\frac{(2 y)^x}{x!} e^{-3 y} & \text { if } y>0 \text { and
  } x=0,1, \ldots, \\ 0 & \text { otherwise. }\end{cases}
$$


a. Verify that $f$ is a joint pmf/pdf.
  (Hint: First, sum over the $x$ values using the well-known formula for the power series expansion of $e^{2 y}$.)
b. Find $\mathbb{P}(X=0)$.

%</lec9-q8>


% ========== Lecture 10 ========== %
%<*lec10-q1>

Suppose the joint pdf of $(X,Y)$ is
  $f(x,y) = c x^2y \cdot \mathbf{1}_{ x^2 \le y <1}$. Find the marginal pdf
  $f(y)$ of $Y$.

%</lec10-q1>


%<*lec10-q2>

Suppose that in a certain drug the concentration of a
particular chemical is a random variable with a continuous
distribution for which the pdf $g(\cdot)$ is defined as follows:

$$
g(x)=
    \begin{cases}
      \frac{3}{8}x^2 & \text{for}\ 0 \leq x \leq 2,\\
      0 & \text{otherwise.}
    \end{cases}
$$

Suppose that the concentrations \(X\) and \(Y\) of the chemical
in two separate batches of the drug are independent random
variables for each of which the pdf is $g(\cdot)$. Determine:

a. The joint pdf of \(X\) and \(Y\).
b. $\mathbb{P}(X=Y)$.
c. $\mathbb{P}(X>Y)$.
d. $\mathbb{P}(X+Y \leq 1)$.

%</lec10-q2>


%<*lec10-q3>

Suppose that $X$ and $Y$ have a discrete joint distribution
for which the joint pmf is defined as follows:

$$
f(x,y) =
    \begin{cases}
      \frac{1}{30}(x+y) & \text{for }\ x=0,2 \text{and}\ y=0,1,2,3,4.\\
      0 & \text{otherwise}
    \end{cases}
$$


a. Determine the marginal pmfs of \( X \) and \( Y \).
b. Are \( X \) and \( Y \) independent?

%</lec10-q3>


%<*lec10-q4>

Imagine that the joint pdf of \(X \) and \(Y\) is as follows:

$$
f(x,y) =
    \begin{cases}
      4xe^{-y} & \text{for }\ 0 \leq x \leq \sqrt{\frac{1}{2}} \text{and}\ 0 \leq y \leq \infty.\\
      0 & \text{otherwise}
    \end{cases}
$$

Are \(X\) and \(Y\) independent?

%</lec10-q4>


%<*lec10-q5>

Take a point \( (X, Y )\) picked arbitrarily from circle \(S\), characterized by:
$$S=\{(x,y):x^2+y^2 \leq 1\}$$
If $f(x,y) = kx^2y^2$, verify that \(X\) and \(Y\) have the same marginal pdf:

$$
f(x) =
    \begin{cases}
      \frac{2kx^2(1-x^2)^{\frac{3}{2}}}{3} & \text{if }\ -1 \leq x \leq 1.\\
      0 & \text{otherwise}
    \end{cases}
$$

%</lec10-q5>


%<*lec10-q6>

Suppose that a point $(X,Y)$ is chosen at random from the following disk
  $S$:
$$S=\{(x,y):(x-1)^2+(y+2)^2 \leq 16\}.$$
Determine:

a. the conditional pdf of $Y$ for every given value of $X$.
b. $\mathbb{P}(Y>0 \mid X=3)$.

%</lec10-q6>


%<*lec10-q7>

Let $Y$ be the rate (calls per hour) at which calls arrive at a
  switchboard. Let $X$ be the number of calls during a two-hour period. Suppose
  that the marginal pdf of $Y$ is

$$
f_2(y)= \begin{cases}e^{-y} & \text { if } y>0 \\ 0 & \text { otherwise }\end{cases}
$$

and that the conditional pmf of $X$ given $Y=y$ is

$$
g_1(x \mid y)= \begin{cases}\frac{(2 y)^x}{x!} e^{-2 y} & \text { if } x=0,1, \ldots, \\ 0 & \text { otherwise. }\end{cases}
$$


(Hint: this is a *mixed* joint distribution.)


a. Find the marginal pmf of $X$. (You may use the formula
  $\int_0^{\infty} y^k e^{-y} d y=k!$)

b. Find the conditional pdf $g_2(y \mid 0)$ of $Y$ given $X=0$.
c. Find the conditional pdf $g_2(y \mid 1)$ of $Y$ given $X=1$.
d. For what values of $y$ is $g_2(y \mid 1)>g_2(y \mid 0)$ ? Does this
  agree with the intuition that the more calls you see, the higher you should
  think the rate is?

%</lec10-q7>


%<*lec10-q8>

Suppose $f(x,y) = e^{-x} \cdot \mathbf{1}_{0\le y \le x <\infty}$. Find the conditional pdf $f_{Y|X}(y|x)$.

%</lec10-q8>


%<*lec10-q9>

Suppose $X \sim \text{Beta}(a,b)$ for constants $a,b >0$, and
  $Y| X=x \sim \text{Binomial}(N,x)$ where $N$ is some fixed constant.
  Find the conditional pdf of $X|Y=y$ for some $y\in 0,...,N$. (Hint: you've
  already seen a simpler version of this in lecture.)

%</lec10-q9>


% ========== Lecture 11 ========== %
%<*lec11-q1>

Suppose that $X$ is a random variable for which the pdf is $f$ and that
  $Y=a X+b$ (for $(a \neq 0)$). Then show that pdf of $Y$ is
  $$
  g(y)=\frac{1}{|a|} f\left(\frac{y-b}{a}\right) \quad \text { for }-\infty<y<\infty,
  $$
and 0 otherwise.

%</lec11-q1>


%<*lec11-q2>

Suppose that a random variable $X$ can have each of the seven values
  $-3,-2,-1,0,1,2,3$ with equal probability.  Determine the pmf of:
  $$Y= X^2-X.$$

%</lec11-q2>


%<*lec11-q3>

Let $X$ be the rate at which customers are served in a queue, following an
  exponential distribution with parameter $\lambda$.

a. Find the pdf $f_Y(y)$ of the average waiting time $Y = 1/X$.
b. Calculate $\int y f_Y(y) dy$. (this is called the "expectation"
  or "expected value" of a random variable).


%</lec11-q3>


%<*lec11-q4>

Let's say a person calls a call center. Let $X$ be the amount of time
  spent on hold and $Y$ be the total duration of the call. Suppose the
  joint density is $f_{X,Y} (x,y)= e^{-y}$ in the region $0 < x <y <\infty$.

a. What is the marginal density of $X$? of $Y$?
b. The amount of time the consultant spends with the caller after the
  call is answered is $W= Y-X$.  What is the probability density of $W$?
  (Hint: use )


%</lec11-q4>


%<*lec11-q5>

Let's say a person calls a call center. Let $X$ be the amount of time
   spent on hold and $Y$ be the total duration of the call. Suppose the
   joint density is $f_{X,Y} (x,y)= e^{-y}$ in the region $0 < x <y <\infty$.

a. What is the marginal density of $X$? of $Y$?
b. The amount of time the consultant spends with the caller after the
  call is answered is $W= Y-X$.  What is the probability density of $W$?

%</lec11-q5>


%<*lec11-q6>

Let random variable $Z$ have a Standard Normal distribution
  $Z \sim \mathcal{N}(0, 1)$ with
  $$f_Z(z) = \frac{1}{\sqrt{2\pi}}
  e^{-z^2/2}.$$

a. Find the probability density of $X = aZ + b$ for constants $a$ and $b$.
b. Find the probability density of $X = Z^2$.

%</lec11-q6>


%<*lec11-q7>

Use the result from the previous problem to prove that the square of a Normal random
  variable is (1) a Chi-squared distribution with 1 degree of freedom, and
  equivalently (2) Gamma($1/2, 1/2$).


- Suppose we have a random variable $U$ that is uniformly distributed on
  the interval $(0, 1)$:
  $$U \sim \text{Uniform}(0, 1).$$

  Find the distribution of $Y$, defined as:
  $$Y = -\frac{1}{\lambda} \ln(U). $$

  (where $\lambda > 0$ is a constant).

%</lec11-q7>


% ========== Lecture 12 ========== %
%<*lec12-q1>

Suppose that $X$ has the uniform distribution on the interval $[0,
  1]$. Construct a random variable $Y= r(X)$ for which the pdf will be
  $$g(y)=
  \begin{cases}
    \frac{3}{8}y^2 &\text{ for } 0 < y < 2\\
    0 & \text{ otherwise.}
  \end{cases}
  $$

%</lec12-q1>


%<*lec12-q2>


a. Suppose that the random variable $X$ is uniform in $[0,1]$.
  Construct a random variable $Y=r_1(X)$ for which the pdf will be:
  $$g(y) =
  \begin{cases} \frac{3}{8} y^2 & \text{ for } 0 < y < 2,\\
    0 & \text{ otherwise}. \end{cases}
b.
  Next, let $Z$ be a random variable whose pdf is:
  $$f(z) =
  \begin{cases} \frac{1}{2} z & \text{ for } 0 < z < 2,\\
  0 & \text{ otherwise}. \end{cases} Construct a random variable
  $W = r_2(Z)$ whose pdf is $g$ from problem (a)? (Hint: start with a
  transformation of $Z$ that makes it distributed as $\text{Unif}(0,1)$, then
  use (a).)

%</lec12-q2>


%<*lec12-q3>

Explain how to use a uniform pseudo-random number generator to generate five
  independent values from the distribution for which the pdf is:

  $$g(y) = \begin{cases}
             (2y+1)/2 & \text{ for } 0 < y < 1,\\
             0 & \text{ otherwise }.
           \end{cases}$$
  (Optional but recommended: try your algorithm on a computer and show the results!
  Submit your code and output for extra credit.)


%</lec12-q3>


%<*lec12-q4>

Suppose $X_1$ and $X_2$ are independent exponential random variables with
  parameter $1$. Find the pdf of $Y = X_1 - X_2$.


%</lec12-q4>


%<*lec12-q5>

Suppose the $X_1, \cdots, X_n$ are random variables that are uniformly
  distributed on $[0,1]$. Derive the pdf of the second largest of them, denoted
  $X_{(n-1)}$.
  (Hint: the event $\{X_{(n-1)} \leq t\}$ is the same as
  $\{\text{At least } n-1 \text{ observations are smaller than or equal to }
  t\}$. Can you see why?)


%</lec12-q5>


%<*lec12-q6>

Show that if $X_1\cdots,X_n$ are independent random variables,


%</lec12-q6>


%<*lec12-q7>

Suppose $X_1, \cdots X_n$ are Unif(0,1) random variables. Find the
  smallest value of the *sample size* $n$ that ensures a high probability ($0.95$ or
  higher) that the maximum $X_{(n)}$ of those $n$ random variables exceeds
  $0.98$.

%</lec12-q7>


% ========== Lecture 13 ========== %
%<*lec13-q1>

Suppose $Z_1,Z_2,...,Z_n$ is a sequence of *independent* random
   variables, and $Z_k \sim \text{Exp}(2^k)$ for $k=1,...,n$.  Find the variance of the
   sample mean $\bar Z$.
   (Hint: use the fact that $\sum_{k=0}^n (1/2)^k = \frac{1/2}{1-(1/2)^{k-1}}$.)

%</lec13-q1>

%<*lec13-q2>

Let $X$ be a random variable that can take only the values $0, 1, 2, \cdots$. Then
 prove that: $$\mathbb{E}[X]= \sum_{n=1}^\infty \mathbb{P}(X \ge n).$$

%</lec13-q2>


%<*lec13-q3>

Suppose that $X$ and $Y$ are independent random variables whose variances
  exist (i.e., are finite) and such that $\mathbb{E}[X]= \mathbb{E}[Y]$.  Show that
  $\mathbb{E}[(X-Y)^2]= V[X] + V[Y]$.

%</lec13-q3>


%<*lec13-q4>

Suppose that $n$ electronic products have been produced; the items are
  independent, and the length of life of each item has the exponential
  distribution with parameter $\lambda$. Determine the expected length of time
  until three items have failed. (Hint: model $Y_i$ to be the time until the
  $i$'th item fails.)

%</lec13-q4>


%<*lec13-q5>

Suppose that a fair coin is tossed repeatedly until a
head is obtained for the first time.

a. What is the expected number of tosses that will be required? (Hint:
  first prove that the probability it will take $n$ tosses or more is
  $(1-p)^{n-1}$. Then, use the result from the in-class exercise.)
b. What is the expected number of tails that will be obtained before the
  first head is obtained?

%</lec13-q5>


%<*lec13-q6>

Suppose that appliances manufactured by a particular company fail at a
  rate of $X$ per year, where $X$ is currently unknown and hence is a random
  variable. If we are interested in predicting how long such an appliance will
  last before failure, we might use the mean of $\frac{1}{X}$. Calculate the mean of
  $Y= \frac{1}{X}$, if the pdf of $X$ is
  $$f(x)= 3x^{2} \cdot \mathbf{1}_{0 < x < 1},$$
  in two ways:

a. First derive the distribution of $Y$, then obtain its expectation.
b. Use the Law of the Unconscious Statistician (LOTUS).

%</lec13-q6>


%<*lec13-q7>

Suppose that a jar contains 10 yellow balls and 15 blue balls, and suppose
  that eight balls are to be selected at random from the jar with replacement.
  How many more yellow balls than blue balls would you expect to
  see?


%</lec13-q7>


%<*lec13-q8>

Derive the variance of a Poisson random variable with parameter
  $\lambda$. (Hint: directly calculate and use $\mathbb{E}[X(X-1)]$.)

%</lec13-q8>


%<*lec13-q9>

Let $X$ be a gamma distributed random variable with parameters $\alpha$
  and $\beta$.

a. Prove that, for $k=1,2,\cdots$, $\mathbb{E}[X^k] = \frac{\Gamma(\alpha +
  k)}{\beta^k \Gamma(\alpha)}$.
b. Use this result to prove that $\mathbb{E}[X] = \alpha/\beta$, and
  $V[X] = \alpha/\beta^2$.
c. What is the mean and variance of an exponential random variable
  with parameter $\beta>0$? (Hint: an exponential distribution is a
  special case of a gamma distribution.)

%</lec13-q9>


%<*lec13-q10>

Suppose that a box contains red balls and blue balls and that the
  proportion of red balls in the box is $p$, ($0 \le p \le 1$). Suppose that $n$
  balls are selected from the box at random

a. If the balls are selected with replacement, what is the expected
  number of red balls?
b. If the balls are selected *without* replacement, what is the
  expected number of red balls?
c. (Extra credit) Suppose now that we sample $n>1$ balls with
  replacement, and let $X$ be the number of red balls in the sample. Then, we
  sample $n$ balls without replacement, and we let $Y$ be the number of red
  balls in that sample.  Prove that $\mathbb{P}(X=n) > \mathbb{P}(Y= n)$.

%</lec13-q10>

% ========== Lecture 14 ========== %
%<*lec14-q1>

Take joint pdf $f_{X,Y}(x,y) = (x+y) \cdot \mathbf{1}_{(x,y)\in [0,1]^2}$ of a
  random vector $(X,Y)$. Derive $\mathbb{E}[Y|X]$. Show all steps.

%</lec14-q1>


%<*lec14-q2>

Suppose that one word is to be selected at random from the sentence “the
  girl put on her beautiful red hat”. If $X$ denotes the number of letters in
  the word that is selected,

a. What is the expected value $X$?
b. Suppose that it was revealed to you that the chosen word was at
  most 4 letters long. What is the expected value now? (In other words, what
  is the conditional expectation?)


%</lec14-q2>


%<*lec14-q3>

Suppose $\theta \sim \text{Unif}(0,2\pi)$, $X=\cos \theta$ and $Y=\sin \theta$.
  Verify that $\operatorname{Cov}(X,Y)=0$.  Is $X\perp Y$?

  (This is a clear example of the two variables $X$ and $Y$ not having a linear
  relationship.)

%</lec14-q3>


%<*lec14-q4>

Suppose that $X$ has the uniform distribution on the interval $[-2, 2]$, and


  $Y = X^6$. Show that $X$ and $Y$ are uncorrelated but not independent
  (where `uncorrelated' means that the correlation is zero).


%</lec14-q4>


%<*lec14-q5>

Suppose that $X$, $Y$, and $Z$ are three random variables such that
  $V[X] = 1$, $V[Y] = 4$, $V[Z] = 8$, $\operatorname{Cov}(X, Y) = 1$, $\operatorname{Cov}(X, Z) = -1$, and
  $\operatorname{Cov}(Y, Z) = 2$.
  Determine $V(3X+2)$ and $V(3X-Y-2Z+1)$.

%</lec14-q5>


%<*lec14-q6>

The joint probability density function of $X, Y$ is
$$
f(x, y)=\frac{1}{y} e^{-(y+x / y)}, \quad 0<x, y<\infty
$$

a. Verify that this is a valid joint probability density function.
(Hint: you might find it useful that $\int_0^{\infty} \lambda e^{-\lambda x} d
x=1$, for any value of $\lambda$.)
b. Find $\operatorname{Cov}(X, Y)$.
c. Find $\mathbb{E}[X|Y]$.
d. **(Optional, extra credit)** Find $V[X|Y]$, and evaluate it at $Y = 5$.

%</lec14-q6>


%<*lec14-q7>

Suppose that a point is chosen at random on a stick of unit length (that
   is, the length of the stick is 1) and that the stick is broken into two pieces
   at that point. Find the expected value of the length of the longer piece.

%</lec14-q7>


%<*lec14-q8>

Suppose that a stick having a length of 3 feet is broken into two pieces,
  and that the point $X$ at which the stick is broken is chosen in accordance with
  pdf $f_X(x)$. What is the correlation between the length of the longer piece and
  the length of the shorter piece? (Hint: the answer is the same no matter what
  $f_X$ is!)

%</lec14-q8>


%<*lec14-q9>

Say we will flip the same coin $n$ times, whose outcomes are
  $X_i|P=p \overset{iid}{\sim} \text{Bernoulli}(p)$. Further suppose that $P \sim \text{Unif}(0,1)$,
  (Hint: what is the distribution of $\sum_{i=1}^n X_i$? Call this a random
  variable $X$.)

a. Compute $\mathbb{E}[\sum \nolimits_{i=1}^n X_i]$, and
b. Compute  $\mathbb{E}[P| \sum \nolimits_{i=1}^n X_i ]$.

%</lec14-q9>


%<*lec14-q10>

Suppose the standard deviations of random variables $X_1,\cdots, X_n$ are
  $\sigma_{X_i}=1$ for $i=1,..,n$, and also, $\rho(X_i,X_j) = 1/4$ for all pairs
  $i\neq j$. Compute the variance of the sum $\sum \nolimits_{i=1}^n X_i $ and
  average $\bar X = \frac{1}{n}\sum_{i=1}^n X_i$.

%</lec14-q10>


%<*lec14-q11>

Suppose $X \sim \text{Exp}(\lambda)$ for some $\lambda >0$, and given $X=x$,
  $Y\sim \text{Pois}(x)$. Find the conditional expectation $\mathbb{E}[X|Y=y]$ for some $y\in\mathbb{Z}_+$.

%</lec14-q11>


%<*lec14-q12>

Suppose $X \sim \text{Beta}(a,b)$ for constants $a,b >0$, and
  $Y| X=x \sim \text{Binom}(N,x)$ where $N$ is some fixed constant.

a. Derive the conditional probability density $f_{X|Y}(x|y)$.
b. Find the conditional expectation $\mathbb{E}[X|Y=y]$ for some
  $y\in \{0,...,N\}$.

%</lec14-q12>


%<*lec14-q13>

Suppose that a point
  $X_1$ is chosen from the uniform distribution on the interval $[0,
  1]$, and that after the value $X_1 = x_1$ is observed, a point
  $X_2$ is chosen from a uniform distribution on the interval $[x_1,
  1]$. Suppose further that additional variables $X_3, X_4, . .
  .$ are generated in the same way. In general, for $j = 1, 2, . .
  .$ , after the value $X_j=x_j$ has been observed,
  $X_{j+1}$ is chosen from a uniform distribution on the interval $[x_j ,
  1]$. Find the value of $\mathbb{E}[X_n]$.

%</lec14-q13>

%<*lec14-q14>

Suppose $\mathbb{E}[Y|X] = aX+b$ for some constants $a,b$. Compute $\mathbb{E}[XY]$.


   Next, suppose $a=0$. What is $\mathbb{E}[Y]$? What is $\operatorname{Cov}(X,Y)$?


%</lec14-q14>

%<*lec14-q15>

Consider a clinical trial in which a number of patients will be treated
   and each patient will have one of two possible outcomes: success or
   failure. Let $P$ be the proportion of successes in $40$ patients, and let
   $X_i = 1$ if the $i$'th patient is a success and $X_i = 0$ if not. Assume that
   the random variables $X_1, X_2, \cdots$ are conditionally independent given
   $P$ with $P(X_i = 1|P)= P$. Let $X = \sum_{i=1}^{40} X_i$, which is the number
   of patients who are successes.

a. What is a good prediction $d$ before the clinical trial starts?
b. What is the best prediction $d(X)$ once the clinical trial has
  ended and you have observed $X$?
c. Let's say that $X=18$ patients were successful. Sketch the (1) pdf
  of $P$ *prior to* the clinical trial, and (2) conditional distribution
  of $P$ having conducted the clinical trial to observe $X=18$.
d. Compute the MSEs of the predictions $d$ and $d(X)$ from (a) and (b)
  respectively, which are $\mathbb{E}[(P-d)^2]$ and
  $\mathbb{E}[(P-d(X))^2|X=18]$. Which is bigger?  Interpret this result in the
  context of the clinical trial.

%</lec14-q15>

% ========== Lecture 15 ========== %
%<*lec15-q1>

What is the moment generating function for a $\text{Unif}(0,1)$ random
  variable? Use it to prove that the expected value of this random variable is
  $1/2$.

%</lec15-q1>


%<*lec15-q2>

What is the moment generating function for a $\text{Unif}(a,b)$ random
  variable? Using this,

a. Prove that the mean (i.e., expectation) of a $\text{Unif}(a,b)$
  random variable is $\frac{a+b}{2}$.
b. Prove that the standard deviation of a $\text{Unif}(a,b)$ is $\frac{(b-a)^2}{12}$.

%</lec15-q2>


%<*lec15-q3>

Suppose that the random variables $X$ and $Y$ are iid and that the mgf of
  each is:
  $$
  \psi(t)=e^{t^2+3 t} \quad \text { for }-\infty<t<\infty .
  $$
  Find the mgf of $Z=2 X-3 Y+4$.

%</lec15-q3>


%<*lec15-q4>

Suppose that $X$ and $Y$ are iid random variable for which the moment generating
  function (mgf) is as follows:
$$\psi(t)=e^{t^2+4t} \ \text{for} \ -\infty < t < \infty $$

a. Find the mean and the variance of $X$.
b. Find the mgf of $Z=3X-2Y+6$.

%</lec15-q4>


%<*lec15-q5>

Take a Poisson random variable $X$ with parameter $\lambda$.


a. What is the mgf of $X$? (Hint: use $\sum_{x=0}^\infty \frac{a^x}{x!} = e^a$.)
b. Prove that both the mean and variance of $X$ are $\lambda$.

%</lec15-q5>


%<*lec15-q6>

Show that the expectation of the sum of two Poisson random variables
  $X\sim \text{Pois}(\lambda_1)$ and $Y \sim \text{Pois}(\lambda_2)$ is
  $\lambda_1 + \lambda_2$, in at least two different ways.

%</lec15-q6>


%<*lec15-q7>

Suppose that $X$ is a random variable for which the mgf is as follows:
  $$
  \psi(t)=\frac{1}{4}\left(3 e^t+e^{-t}\right) \quad \text { for
  }-\infty<t<\infty .
  $$
  Find the mean and the variance of $X$.

%</lec15-q7>


%<*lec15-q8>

Suppose it is known that $30$ percent of the items in a large manufactured
  lot are of poor quality. Suppose also that a random sample of $n$ items is to
  be taken from the lot, and let $Q_n$ denote the proportion of the items in the
  sample that are of poor quality.

  Find a value of $n$ such that
  $\mathbb{P}(0.2 \le Q_n \le 0.4) \ge 0.75$ by using:

a. the Chebyshev inequality and
b. the exact probabilities of the binomial distribution. (Hint: you may
  use R code and the functions \texttt{pbinom()} and \texttt{dbinom()}
  useful. Otherwise, you may choose to use a calculator.)
c. Which of (a) and (b) tell you a larger $n$? What is your
  interpretation of this result?

%</lec15-q8>


%<*lec15-q9>

Let $X$ be the number of calls arriving at a call center in the next
  hour. $X$ is distributed as a Poisson random variable with mean parameter $3$. Use
  Chebyshev's inequality to show that the probability that there will be more
  than $6$ calls is smaller than $1/3$.


%</lec15-q9>


% ========== Lecture 16 ========== %
%<*lec16-q1>

Take $50$ random variables $X_1, \cdots, X_n$, each i.i.d. with mean $\mu$
  and variance $\sigma^2=9$. Find the 90\% probability region of $\bar X - \mu$,
  using (a) Chebyshev's inequality, and (b) Central limit theorem. Which is
  wider, and provide your interpretation: why is it wider?


%</lec16-q1>


%<*lec16-q2>

Booklets at a factory will be packaged into boxes of 100. On average, the
  booklets weigh 1 ounce, with a standard deviation of $0.05$ ounce. The
  manufacturer is interested in calculating:

  $$\mathbb{P}(100 \text{ booklets weigh more than } 100.4 \text{ ounces}),$$

  a number that would help the manufacturer determine whether too many booklets
  are being put into a box. Explain how you would calculate the (approximate)
  value of this probability. (Mention any relevant theorems or assumptions
  needed.)

%</lec16-q2>


%<*lec16-q3>

Find the *limiting distribution* of $\bar X_n$, for i.i.d. random
  variables $X_1, \cdots, X_n$ that are distributed as $\text{Unif}(a,b)$.


%</lec16-q3>

%<*lec16-q4>

Suppose that the height of men has mean 68 inches and standard deviation
  2.6 inches. We draw 100 men at random. Find (approximately) the probability
  that the average height of men in our sample will be at least 68 inches.\\

  (Hint: you might want to study how to calculate normal probabilities. See the
  textbook -- e.g. section 5 and Example 5.6.4)

%</lec16-q4>
%<*lec16-q5>

%</lec16-q5>
%<*lec16-q6>

%</lec16-q6>
%<*lec16-q7>

%</lec16-q7>
%<*lec16-q8>

%</lec16-q8>
%<*lec16-q9>

%</lec16-q9>
%<*lec16-q10>

%</lec16-q10>

% ========== Lecture 17 ========== %
%<*lec17-q1>

%</lec17-q1>
%<*lec17-q2>

%</lec17-q2>
%<*lec17-q3>

%</lec17-q3>
%<*lec17-q4>

%</lec17-q4>
%<*lec17-q5>

%</lec17-q5>
%<*lec17-q6>

%</lec17-q6>
%<*lec17-q7>

%</lec17-q7>
%<*lec17-q8>

%</lec17-q8>
%<*lec17-q9>

%</lec17-q9>
%<*lec17-q10>

%</lec17-q10>

% ========== Lecture 18 ========== %
%<*lec18-q1>

%</lec18-q1>
%<*lec18-q2>

%</lec18-q2>
%<*lec18-q3>

%</lec18-q3>
%<*lec18-q4>

%</lec18-q4>
%<*lec18-q5>

%</lec18-q5>
%<*lec18-q6>

%</lec18-q6>
%<*lec18-q7>

%</lec18-q7>
%<*lec18-q8>

%</lec18-q8>
%<*lec18-q9>

%</lec18-q9>
%<*lec18-q10>

%</lec18-q10>

% ========== Lecture 19 ========== %
%<*lec19-q1>

%</lec19-q1>
%<*lec19-q2>

%</lec19-q2>
%<*lec19-q3>

%</lec19-q3>
%<*lec19-q4>

%</lec19-q4>
%<*lec19-q5>

%</lec19-q5>
%<*lec19-q6>

%</lec19-q6>
%<*lec19-q7>

%</lec19-q7>
%<*lec19-q8>

%</lec19-q8>
%<*lec19-q9>

%</lec19-q9>
%<*lec19-q10>

%</lec19-q10>
